\begin{figure*}[t]
\centering
\scriptsize
\setlength{\fboxsep}{1em}
\setlength{\fboxrule}{1pt}
\fcolorbox{black}{white}{
\parbox{0.955\textwidth}{
\begin{alltt}


You will be provided instructions for completing a task followed by a task to complete.

Instructions:

1. Understand the question: The questions in this dataset involve arithmetic operations and real-life scenarios. Make sure you understand what operation is being asked for (addition, subtraction, multiplication, division) and the context of the question.

2. Perform the operation: Once you understand the question, perform the necessary arithmetic operation. This could involve simple calculations or more complex ones depending on the question. 

3. Check your work: Always double-check your calculations to ensure accuracy. 

4. Formulate the answer: The answer should be a string that represents the final result of the calculation. Make sure your answer is clear and concise.

5. Topics to know: Basic arithmetic operations (addition, subtraction, multiplication, division), understanding of real-life scenarios (counting items, calculating money, etc.)

6. Steps to answer the questions: Read the question carefully, identify the arithmetic operation and the context, perform the operation, check your work, and write the answer as a string.

Question: For Halloween Debby and her sister combined the candy they received. Debby had 32 pieces of candy while her sister had 42. If they ate 35 pieces the first night, how many pieces do they have left?

Follow the instructions to generate an explanation that reasons towards the correct answer to the task above. End the explanation with the correct answer. [/INST] Explanation: 

To solve this problem, we need to use subtraction. Debby had 32 pieces of candy, and her sister had 42 pieces, so together they had 32 + 42 = 74 pieces of candy.

They ate 35 pieces the first night, so the amount of candy they have left is 74 - 35 = 39 pieces.

Therefore, the answer is 39 pieces of candy.

Answer: 39 pieces of candy.

Therefore, the answer to the task is below. Give the answer in the shortest form possible that will still be correct.

Answer: 39
\end{alltt}
}}
\caption{A complete running example of MultiArith.}
\label{fig:multi_arith_example}
\end{figure*}
